%=========== ΛΥΣΗ ΘΕΜΑΤΟΣ Γ.96 ===========
\begin{thema}{Γ}

\begin{center}
\begin{tikzpicture}[scale=0.8]
  \draw[help lines, gray!30] (-1,-1) grid (5,4);
  \draw[-latex] (-1.3,0) -- (5.3,0) node[right] {$x$};
  \draw[-latex] (0,-1.3) -- (0,4.3) node[above] {$y$};
  \foreach \x in {1,2,3,4} \draw (\x,0.1) -- (\x,-0.1) node[below,font=\footnotesize] {$\x$};
  \foreach \y in {1,2,3} \draw (0.1,\y) -- (-0.1,\y) node[left,font=\footnotesize] {$\y$};
  
  \draw[dashed, gray!50] (-1,-1) -- (4,4) node[right, font=\footnotesize, fill=white, inner sep=1pt] {$y=x$};
  
  \fill[cyan!10, opacity=0.5, smooth, tension=0.4]
    (0,0) -- plot coordinates {(0,0) (1,1.5) (2,2.2) (3,2.7) (4,3)} -- (4,0) -- cycle;
    
  \fill[yellow!20, opacity=0.4, smooth, tension=0.4]
    (0,0) -- plot coordinates {(0,0) (1,1.5) (2,2.2) (3,2.7) (4,3)} -- (0,3) -- cycle;
  
  \draw[very thick, blue!70!black, smooth, tension=0.4]
    plot coordinates {(0,0) (1,1.5) (2,2.2) (3,2.7) (4,3)};
    
  \node[blue!70!black, fill=white, inner sep=1pt] at (3,1) {$E_1 = \int_0^4 f(x) dx = 4$};
  \node[orange!80!black, fill=white, inner sep=1pt] at (1.5,2.4) {$E_2 = \int_0^3 f^{-1}(y) dy = 8$};
  
  \draw[dashed, purple!60!black, very thick] (4,0) -- (4,3) -- (0,3);
  \node[purple!60!black, fill=white, inner sep=1pt] at (2, 3.6) {$\text{Ορθογώνιο } 4 \times 3 \ (\text{Εμβαδόν } 12)$};
  
  \draw[dashed, red!70!black, thick] (0,0) -- (4,3) node[pos=0.4, below right, font=\footnotesize, fill=white, inner sep=1pt] {$\text{Χορδή κάτω από την κοίλη } C_f$};
  
  \filldraw[blue!70!black] (0,0) circle (2pt);
  \filldraw[blue!70!black] (4,3) circle (2pt) node[below right,font=\footnotesize] {$(4,3)$};
\end{tikzpicture}
\end{center}

\begin{erwthma}

\item Η συνάρτηση $f$ δίνεται ρητά ως \textbf{γνησίως αύξουσα} στο πεδίο ορισμού της $[0, 4]$. Λόγω της γνήσιας μονοτονίας, η $f$ είναι συνάρτηση 1-1, γεγονός που εγγυάται ότι \textbf{αντιστρέφεται}.
- Το πεδίο ορισμού της αντίστροφης συνάρτησης $f^{-1}$ ταυτίζεται με το σύνολο τιμών της $f$.
- Εφόσον η $f$ είναι συνεχής και γνησίως αύξουσα στο $[0, 4]$, το σύνολο τιμών της είναι το $[f(0), f(4)] = [0, 3]$. 
Άρα το πεδίο ορισμού της $f^{-1}$ είναι το \textbf{$[0, 3]$}.

\item Δίνεται η βασική γεωμετρική/αλγεβρική ιδιότητα που συνδέει τα εμβαδά μιας συνάρτησης 1-1 με την αντίστροφή της στο περιβάλλον ορθογωνίου:
\[ \int_0^4 f(x)\,\d x + \int_{f(0)}^{f(4)} f^{-1}(y)\,\d y = 4 \cdot f(4) - 0 \cdot f(0) \]
Αντικαθιστώντας τα γνωστά δεδομένα $f(0)=0$, $f(4)=3$, και το δεδομένο $\int_0^4 f(x)\,\d x = 4$:
\begin{align*}
4 + \int_0^3 f^{-1}(y)\,\d y &= 4 \cdot 3 \\
4 + \int_0^3 f^{-1}(y)\,\d y &= 12 \\
\int_0^3 f^{-1}(y)\,\d y &= 12 - 4 = \mathbf{8 \ \text{τ.μ.}}
\end{align*}

\item Η \textbf{μέση τιμή} μιας συνεχούς συνάρτησης σε ένα κλειστό διάστημα $[a, b]$ δίνεται από τον τύπο $\overline{f} = \frac{1}{b-a} \int_a^b f(x)\,\d x$.
Για την $f$ στο $[0, 4]$:
\[ \overline{f} = \frac{1}{4 - 0} \int_0^4 f(x)\,\d x = \frac{1}{4} \cdot 4 = \mathbf{1} \]
Σύμφωνα με το \textbf{Θεώρημα Μέσης Τιμής του Ολοκληρωτικού Λογισμού} (Θ.Ε.Τ. / Θ.Μ.Τ.Ο.Λ.), εφόσον η $f$ είναι συνεχής στο $[0, 4]$, υπάρχει τουλάχιστον ένα σημείο $\xi \in (0, 4)$ τέτοιο ώστε η συνάρτηση να λαμβάνει τη μέση τιμή της, δηλαδή:
\[ \mathbf{f(\xi) = 1} \]

\item Εξετάζουμε τη σχετική θέση της καμπύλης με τη χορδή (το ευθύγραμμο τμήμα) που ενώνει τα άκρα $(0,0)$ και $(4,3)$. 
Από την εκφώνηση γνωρίζουμε ότι η $C_f$ είναι \textbf{κοίλη} (ή αλλιώς, στραμμένα τα κοίλα προς τα κάτω, $f''(x) < 0$). 
Μια θεμελιώδης γεωμετρική ιδιότητα των κοίλων συναρτήσεων είναι ότι η γραφική τους παράσταση βρίσκεται πάντα \textbf{περισσότερο κυρτωμένη "προς τα πάνω" σε σχέση με οποιαδήποτε χορδή της}. 
Δηλαδή, για κάθε $x \in (0, 4)$ του ενδιάμεσου πεδίου, η συνάρτηση διατηρεί τιμές \textbf{υψηλότερες} από αυτές της αντίστοιχης διαγωνίου (εδώ ευθεία $y = \frac{3}{4}x$).
Αποδεικνύεται λοιπόν άμεσα (και οπτικά) ότι η χορδή $(0,0)-(4,3)$ **περνά και βρίσκεται εξ' ολοκλήρου ΚΑΤΩ** από τη γραφική παράσταση της $C_f$.
\end{erwthma}
\end{thema}
